\documentclass[10pt,a4paper]{amsart}
\usepackage[margin=19mm]{geometry}
\usepackage{amsmath,amssymb,mathtools}
\usepackage[scr=rsfs,cal=euler]{mathalfa}
\usepackage{xfrac}
\usepackage[unicode,hidelinks,bookmarks=false]{hyperref}
\newcommand{\ie}{i.e.\ }
\newcommand{\cf}{cf.\ }
\newcommand{\resp}{respectively\ }
\newcommand{\Subsection}{\subsection}
\providecommand{\nobreakdash}{\nobreak\hbox{-}}
\setlength{\emergencystretch}{2em}
\multlinegap0pt
\usepackage{bm}
\usepackage{bbm}
\usepackage{dsfont}
\usepackage{xspace}
\usepackage{cancel}
\usepackage{mathtools}
\usepackage{amsthm}
\makeatletter
\@ifundefined{proposition}{\newtheorem{proposition}{Proposition}}{}
\makeatother
\title[Magic and Non-Clifford Gates in Topological Quantum Field Th]{Magic and Non-Clifford Gates in Topological Quantum Field Theory}
\author{William Munizzi, Howard J. Schnitzer}
\date{Source: arXiv:2604.14271v1 (CC BY 4.0)}
\begin{document}
\maketitle

\noindent\emph{Source and licence.} Magic and Non-Clifford Gates in Topological Quantum Field Theory. Version arXiv:2604.14271v1 (CC BY 4.0), distributed under the Creative Commons Attribution 4.0 International licence, as recorded in the arXiv licence metadata for this exact version and shown on the abstract page https://arxiv.org/abs/2604.14271v1. Full rights notice: licenses/arxiv-2604.14271-CC-BY-4.0.txt.\par\smallskip

\noindent\textit{Editorial scope.} Counted unit: the Proposition whose source label is \texttt{prop:nonlocal} in the article (source0.tex lines 583--596, source label \texttt{prop:nonlocal}), delivered together with the article's own complete original proof of it (source0.tex lines 598--656, one \texttt{proof} environment of 59 lines). No mathematics is written, repaired or re-derived: statement and proof are the article's own text line for line. Required context retained uncounted: OpSchmidtDecomp (lines 426--428); ELimReduction (lines 432--434); NonLocalMagicCriterion (lines 440--442); GateDef (lines 539--541). Every definition, display and label of those lines is retained unchanged. Omitted from this package and not redistributed: 1--425, 429--431, 435--439, 443--538, 542--582, 597--597, 657--1063. Nothing inside a retained excerpt is omitted or altered: every retained body is delivered exactly as the article's lines are written, so no omission receipt of any kind accompanies this package. Citations: the retained excerpts carry no citation command at all, so no bibliography entry of the article is transcribed and no reference is redistributed. Reference adaptation: none was needed and none was made; every reference inside the delivered unit resolves inside it, so no unresolved reference or citation is produced. Printed slips kept literal: no printed word, symbol, hypothesis or number of the counted statement or of its proof was altered. Layout and class: the article's own class is \texttt{article} with packages including inputenc, fullpage, cite, caption, subcaption, babel, amsmath, physics, amsthm, amsfonts, algorithm, algpseudocode, amssymb, mathrsfs; the wrapper is the lane's pinned \texttt{amsart} editorial wrapper at 10pt on A4, which restates the theorem-like environments the retained text needs and adds the article's own macro definitions that the retained text needs (none), with extra wrapper packages (bm, bbm, dsfont, xspace, cancel, mathtools). The delivered numbering is the unit's own. Assets: no retained span contains a figure environment, a graphics-inclusion command, a PDF-page inclusion, a table or a TikZ picture; no image file is copied into this package, the package declares no asset dependency and no image file is redistributed with it. Licence: version arXiv:2604.14271v1 of this article is distributed under the Creative Commons Attribution 4.0 International licence, as recorded in the arXiv licence metadata for this exact version and shown on the abstract page https://arxiv.org/abs/2604.14271v1. A full rights notice is delivered at licenses/arxiv-2604.14271-CC-BY-4.0.txt. Characters: every retained excerpt is pure ASCII, so the delivered engine, XeLaTeX with its default Latin Modern text font, reports no missing character.
\begin{equation}\label{OpSchmidtDecomp}
    U/\sqrt{d^n} = \sum_i \sqrt{\lambda_i}\, V_i \otimes W_i, 
\end{equation}

\begin{equation}\label{ELimReduction}
    E_{\mathrm{lin}}(U) = 1 - \sum_i \lambda_i^2.  
\end{equation}

\begin{equation}\label{NonLocalMagicCriterion}
E_{\mathrm{lin}}\left(U^\dagger P U \right) = 0 \quad \forall P \in \mathcal{P}_n \iff U = (V \otimes W)\,C, \quad V \in \mathcal{U}_A,\; W \in \mathcal{U}_B,
\end{equation}

\begin{equation}\label{GateDef}
U(\theta) = \exp\left(-i\,\frac{\theta}{2}\,X_1\otimes X_2\right), \qquad \theta \in [0,\pi].
\end{equation}

\begin{proposition}\label{prop:nonlocal}
For $\theta \in (0,\pi/2)\cup(\pi/2,\pi)$, the Ising interaction gate $U(\theta)$ defined in Eq.\ \eqref{GateDef} generates non-local magic. Specifically for $P = Z_1\otimes\mathbb{I}_2$,
%
\begin{equation}\label{ConjugationResult}
U(\theta)^\dagger\,(Z_1\otimes\mathbb{I}_2)\,U(\theta) = \cos(\theta)\,(Z_1\otimes\mathbb{I}_2) - \sin(\theta)\,(Y_1\otimes X_2),
\end{equation}
%
and the operator linear entropy of the conjugated Pauli string satisfies
%
\begin{equation}\label{ElinResult}
E_{\mathrm{lin}}\left(U(\theta)^\dagger\,(Z_1\otimes\mathbb{I}_2)\,U(\theta)\right) = \frac{1}{2}\sin^2(2\theta) > 0.
\end{equation}
%
\end{proposition}

\begin{proof}
Let $G \equiv X_1\otimes X_2$ such that $U(\theta)$ and $U^{\dagger}(\theta)$ may be written
%
\begin{equation}\label{Rewriting}
\begin{split}
    U(\theta) &= \cos(\theta/2)\,\mathbb{I} - i\sin(\theta/2)\,G,\\
    U(\theta)^\dagger &= \cos(\theta/2)\,\mathbb{I} + i\sin(\theta/2)\,G.\\
\end{split}
\end{equation}
%
Conjugating $P = Z_1\otimes\mathbb{I}_2$ by $U(\theta)$ and expanding, we have
%
\begin{align}\label{ConjugationExpansion}
U(\theta)^\dagger(Z_1\otimes\mathbb{I}_2)U(\theta) &= \cos^2\frac{\theta}{2}\,(Z_1\otimes\mathbb{I}_2) + i\cos\frac{\theta}{2}\sin\frac{\theta}{2}\bigl[G,\,Z_1\otimes\mathbb{I}_2\bigr] \notag\\
&\quad \quad + \sin^2\frac{\theta}{2}\,G\,(Z_1\otimes\mathbb{I}_2)\,G.
\end{align}
%
Employing the Pauli relations $X_1 Z_1 = iY_1$ and $Z_1 X_1 = -iY_1$, we can expand the commutator $\bigl[G,\,Z_1\otimes\mathbb{I}_2\bigr]$ as
%
\begin{equation}
\bigl[G,\,Z_1\otimes\mathbb{I}_2\bigr] = (X_1 Z_1 - Z_1 X_1)\otimes X_2 = 2iY_1\otimes X_2,
\end{equation}
%
and therefore the second term in Eq.\ \eqref{ConjugationExpansion} contributes
$-\sin(\theta)\,(Y_1\otimes X_2)$. For the final term in Eq.\ \eqref{ConjugationExpansion} we have
%
\begin{equation}
    G(Z_1\otimes\mathbb{I}_2)G = (X_1 Z_1 X_1)\otimes\mathbb{I}_2 = -Z_1\otimes\mathbb{I}_2.
\end{equation}
%
Applying the relation $\cos^2(\theta/2) - \sin^2(\theta/2) = \cos(\theta)$ then gives Eq.\ \eqref{ConjugationResult}. We can observe that the local operator $Z_1\otimes\mathbb{I}_2$ has been rotated into a superposition with the bipartite operator $Y_1\otimes X_2$, with mixing parameter $\theta$.

Establishing positivity of the operator linear entropy in Eq.\ \eqref{ElinResult}, we first write the conjugated operator%
%
\footnote{Conjugation by a unitary preserves unitarity. The operator $U^\dagger P U$ is unitary, therefore the operator Schmidt decomposition of Eq.\ \eqref{OpSchmidtDecomp} carries the normalization $(U^\dagger P U)/\sqrt{d^n} = (U^\dagger P U)/2$.} %
%
from Eq.\ \eqref{ConjugationResult} in its Schmidt form
%
\begin{equation}\label{SchmidtFormOp}
\frac{1}{2}\left(\cos(\theta)\,(Z_1\otimes\mathbb{I}_2) - \sin(\theta)\,(Y_1\otimes X_2)\right) = \cos(\theta)\,\frac{Z_1}{\sqrt{2}}\otimes\frac{\mathbb{I}_2}{\sqrt{2}} - \sin(\theta)\,\frac{Y_1}{\sqrt{2}}\otimes\frac{X_2}{\sqrt{2}}.
\end{equation}
%
Both terms in Eq.\ \eqref{SchmidtFormOp} have orthogonal Hilbert-Schmidt inner product, since $\mathrm{Tr}(Z^\dagger Y) = 0$ and $\mathrm{Tr}(\mathbb{I}^\dagger X) = 0$, yielding Schmidt eigenvalues 
%
\begin{equation}
\begin{split}
    \mu_1 &= \cos^2(\theta),\\
    \mu_2 &= \sin^2(\theta).
\end{split}
\end{equation}
%
Applying the reduced operator linear entropy form in Eq.\ \eqref{ELimReduction} gives
%
\begin{equation}
E_{\mathrm{lin}}(U(\theta)) = 1 - \mu_1^2 - \mu_2^2 = \frac{1}{2}\sin^2(2\theta),
\end{equation}
%
which is strictly positive for $\theta\in(0,\pi/2)\cup(\pi/2,\pi)$. Therefore the Ising interaction gate $U(\theta)$ generates non-local magic according to the criterion in Eq.\ \eqref{NonLocalMagicCriterion}.
\end{proof}

\end{document}
